\chapter{Configuration and robot profile}
\label{ch:config}

\section{The robot profile}

Every node reads the same \param{robot.*} parameters into an immutable
\code{dharma_utils.robot.RobotProfile}. The product passes them once, in its robot params file,
to every node through the launch argument \code{robot_params}:

\begin{lstlisting}
/**:
  ros__parameters:
    robot:
      name: quadruped_a
      kinematics: omni          # omni: vx, vy, wz; diff: vx, wz
      base_height: 0.35         # m, base_link above the floor when standing
      footprint: {front: 0.40, rear: 0.40, half_width: 0.28}
      limits:
        max_vx: 0.8             # m/s
        max_vy: 0.4             # m/s (0 for diff)
        max_wz: 1.0             # rad/s
        max_accel: 0.5          # m/s^2
        max_decel: 0.6          # m/s^2, braking: measure it on the robot
        max_yaw_accel: 1.5      # rad/s^2, yaw braking
        reaction_time: 0.25     # s, command to braking (sensing + transport)
      frames: {map: map, odom: odom, base: base_link}
\end{lstlisting}

\readonlynote

% Generated by docs/manual/tools/reference.py; do not edit.
\begin{small}
\begin{longtable}{@{}p{0.28\textwidth}p{0.12\textwidth}p{0.08\textwidth}p{0.11\textwidth}p{0.29\textwidth}@{}}
\toprule
Parameter & Default & Unit & Range & Description \\
\midrule
\endhead
\bottomrule
\endfoot
\multicolumn{5}{@{}l}{\textbf{every node}} \\
\code{robot.name} & \code{robot} &  &  & robot model, for logs \\
\code{robot.kinematics} & \code{omni} &  & omni, diff & omni: vx, vy, wz; diff: vx, wz \\
\code{robot.base_height} & 0.35 & m & 0.01\,..\,3 & base\_link above the floor when standing \\
\code{robot.footprint.front} & 0.4 & m & 0.01\,..\,5 & base\_link to the front edge \\
\code{robot.footprint.rear} & 0.4 & m & 0.01\,..\,5 & base\_link to the rear edge \\
\code{robot.footprint.half_width} & 0.28 & m & 0.01\,..\,5 & half of the body width \\
\code{robot.limits.max_vx} & 0.8 & m/s & 0\,..\,10 & forward / backward speed \\
\code{robot.limits.max_vy} & 0.4 & m/s & 0\,..\,10 & sideways speed (omni only) \\
\code{robot.limits.max_wz} & 1 & rad/s & 0\,..\,10 & yaw rate \\
\code{robot.limits.max_accel} & 0.5 & m/s$^{2}$ & 0.01\,..\,20 & linear acceleration \\
\code{robot.limits.max_decel} & 0.6 & m/s$^{2}$ & 0.05\,..\,20 & braking deceleration \\
\code{robot.limits.max_yaw_accel} & 1.5 & rad/s$^{2}$ & 0.05\,..\,50 & yaw braking \\
\code{robot.limits.reaction_time} & 0.25 & s & 0\,..\,2 & command to braking delay (sensing + transport) \\
\code{robot.frames.map} & \code{map} &  &  & global, fixed frame \\
\code{robot.frames.odom} & \code{odom} &  &  & continuous odometry frame \\
\code{robot.frames.base} & \code{base_link} &  &  & robot body frame \\
\end{longtable}
\end{small}


\subsection{Derived quantities}

Nothing derived from the robot is configured separately:
\begin{center}
\begin{tabularx}{\textwidth}{@{}lX@{}}
  \toprule
  Quantity & Derivation and users \\
  \midrule
  inscribed radius & $r_{\mathrm{in}} = \min(f, b, w)$: lethal radius of the planner (plus margin),
    inner edge of the rotation ring \\
  circumscribed radius & $r_{\mathrm{c}} = \sqrt{\max(f,b)^2 + w^2}$: oriented checks of the
    planner, outer edge of the ring swept when turning, creep cap in REMOTE \\
  self filter & footprint grown by a margin: perception, mapping \\
  floor & \frm{base_link} minus \param{robot.base_height}: height bands of perception, mapping,
    scan matching \\
  lateral speed & $v_y = 0$ for \code{diff}: envelope, controllers, approach \\
  braking model & $D(v) = v\,\tr + v^2/(2a)$ with \param{max_decel} and \param{reaction_time}:
    speed profile, path shift, motion guard, protective field \\
  jump limits & \param{max_vx}, \param{max_wz}: odometry monitor, mapping jump gate \\
  \bottomrule
\end{tabularx}
\end{center}

For the default footprint $r_{\mathrm{in}} = 0.28$~m and $r_{\mathrm{c}} = 0.49$~m. The 0.1
planner used a radius of 0.35~m for this footprint (finding N1): a path it accepted could bring
a corner of the body about 14~cm into a wall when the robot turned or moved sideways in a
narrow corridor.

The profile validates cross-field consistency: \param{max_vy} must be 0 for \code{diff}
kinematics and the three frames must differ.

\subsection{Measuring the limits}

\param{max_decel}, \param{max_yaw_accel} and \param{reaction_time} bound every speed limit in
the stack. Measure them on the robot in the worst case (slippery floor, full payload, low
battery): command a step from cruise speed to zero and record the time until deceleration
starts (reaction, including sensing and transport latency) and the stopping distance; use the
smallest deceleration and the largest reaction time observed, then add margin.
\todo{measurement procedure with a script and a recorded example (phase 2)}

\section{Typed parameters}

Each node declares its configuration as frozen dataclasses with
\code{dharma_utils.config.param} fields (default, \code{doc}, \code{unit}, \code{min}, \code{max}, \code{choices}, \code{live})
and \code{nested(...)} sections (dotted names such as \code{speed.cruise}). \code{DharmaNode}:
\begin{itemize}
  \item declares every field as a ROS parameter with a descriptor (description, unit, range,
    read-only for \code{live=False}): \texttt{ros2 param describe <node> <name>} shows it;
  \item validates all values at start (ranges, choices, types and the cross-field
    \code{validate()} of the dataclass) and refuses to start on a safety value that fails;
  \item validates every \texttt{ros2 param set} the same way and rejects an invalid value with
    the reason, keeping the old one;
  \item reports parameters in the YAML files that no dataclass declares (typos).
\end{itemize}

Safety values (margins, timeouts, limits) are read-only at runtime. Defaults live in the
dataclasses; the package's \code{config/<pkg>.yaml} contains only what differs and
documentation. Each subsystem chapter lists its parameters, generated from the dataclasses by
\code{docs/manual/tools/reference.py}.

\section{Parameter files and launch}

A product layers its files over the package defaults:
\begin{enumerate}
  \item the package defaults \code{config/<pkg>.yaml};
  \item the robot profile (\code{robot_params}), shared by every node;
  \item a stack or site file (\code{params_file}) with per-node overrides;
  \item non-empty launch arguments (an empty argument never overrides a file; a file argument
    that is set but missing fails the launch).
\end{enumerate}

Paths accept \code{~} and \code{$VAR}. Paths received through services (tags files, routes,
sites) must resolve inside \code{sites_dir} (\code{dharma_utils.io.resolve_inside}); a request
that escapes it is refused.
